\section{Introduction}
An economy is a large multi-agent system. Households choose how much to work and save, while aggregate choices
determine capital, labour, wages and returns; taxes and redistribution further change incentives. Traditional
macroeconomic models assume agents know the model and solve for equilibrium, whereas learning-based economic
simulators let agents adapt through interaction.

In the Bewley--Huggett--Aiyagari family~\cite{aiyagari,huggett}, households face idiosyncratic income risk, save in a
single asset and take prices as given; equilibrium is obtained through household dynamic programming and a
market-clearing fixed point. In the infinite-population limit this is a mean-field game~\cite{lasry,light}, and
reinforcement learning has been used to find regularized mean-field equilibria~\cite{xu}. Conversely, the AI
Economist~\cite{aiecon} and TaxAI~\cite{taxai} place learning agents inside economic simulators, and Curry et
al.~\cite{curry} study general equilibrium with deep RL. TaxAI in particular benchmarks MARL households and
governments in a Bewley--Aiyagari simulator with up to $10^4$ households, but it has no exact equilibrium to compare
the learned behaviour against.

Our project connects the two views. We build one finite-$N$ economy, solve it with a known-model benchmark (B1), train
decentralized MARL households in the \emph{same} environment, and measure gaps in aggregates, individual policies,
welfare, inequality and exploitability. The central question is:
\emph{how close does the learned economy get to the known-model competitive equilibrium, and how does this depend on
$N$, on the information available to households, and on the tax regime?}

\textbf{Contributions.} (i) An exactly-accounted, vectorized simulator whose resource constraint is verified to
$10^{-9}$ by an automated test; (ii) a dynamic-programming ground truth (EGM, histogram distribution, Brent fixed
point), cross-validated against value-function iteration; (iii) shared-policy IPPO and MAPPO households trained under
four observation levels, five population sizes and seven tax regimes (49 training runs); (iv) aggregate, policy,
consumption-equivalent welfare, inequality and two exploitability metrics (an exact DP best response and a PPO best
response); (v) quantitative tests of four hypotheses H1--H4. All numbers below are produced by the code in our
repository (Section~\ref{sec:repo}).
